% =====================================================================
% RESUMO E ABSTRACT. Todo número aqui aparece no corpo do texto
% (Caps. 5 e 6). Se os resultados mudarem, atualizar os dois.
% =====================================================================

\chapter*{Resumo}

Os modelos de linguagem de grande escala tornaram rápida a produção de questões de Matemática,
mas escrevem com a mesma fluência respostas corretas e incorretas, o que transfere ao professor uma
revisão minuciosa e expõe a sala de aula a erros que não se anunciam. Esta dissertação investiga como
apoiar o professor do Ensino Médio na geração de questões com esses modelos garantindo correção
matemática verificável e adequação didática. Para isso propõe, implementa e avalia um sistema
multiagente em que um modelo de linguagem redige a questão a partir de uma especificação ancorada em
uma habilidade da BNCC, um verificador simbólico baseado no SymPy recalcula o gabarito de forma
independente, um crítico didático a avalia segundo uma rubrica derivada de Pólya, da Teoria das
Situações Didáticas, das taxonomias de Bloom e SOLO e da BNCC, e um orquestrador decide entre
aprovar, revisar ou descartar. O recorte abrange quinze habilidades de funções e progressões. Em um
acervo de 90 especificações, 85 resultaram em questões aprovadas, 72 delas com o gabarito conferido
simbolicamente, a um custo médio de cerca de oito centavos de dólar por questão. Avaliadas por cinco
professores da rede pública, as questões foram aceitas com ou sem ajuste em oito de cada dez
decisões, e as recusas se deveram ao nível da turma e não a erros matemáticos. O único erro
matemático do acervo passou pela verificação porque estava na tradução do enunciado para a
formalização simbólica, o que confirma o fenômeno do erro silencioso e delimita o alcance da
verificação. A partir desse resultado o sistema passou a declarar ao professor com precisão o que foi
conferido em cada questão. O produto educacional reúne o sistema, um banco curado de 85 questões e um
manual do professor, todos de acesso livre.

\medskip
\noindent\textbf{Palavras-chave:} Educação Matemática. Inteligência artificial. Modelos de linguagem.
Sistemas multiagentes. Verificação simbólica. Elaboração de questões.

\chapter*{Abstract}

Large language models have made it fast to produce Mathematics questions, but they write correct and
incorrect answers with the same fluency, which shifts a meticulous review onto the teacher and exposes
the classroom to errors that do not announce themselves. This dissertation investigates how to support
high school teachers in generating questions with these models while ensuring verifiable mathematical
correctness and didactic adequacy. To this end it proposes, implements and evaluates a multi-agent
system in which a language model drafts the question from a specification anchored in a skill of the
Brazilian National Common Curricular Base, a symbolic verifier built on SymPy independently recomputes
the answer key, a didactic critic assesses the question against a rubric derived from Pólya, the
Theory of Didactical Situations, the Bloom and SOLO taxonomies and the curriculum, and an orchestrator
decides whether to approve, revise or discard it. The scope covers fifteen skills on functions and
sequences. In a set of 90 specifications, 85 produced approved questions, 72 of them with a
symbolically checked answer key, at an average cost of about eight US cents per question. Reviewed by
five public school teachers, the questions were accepted with or without changes in eight out of ten
decisions, and rejections were due to the level of the students rather than to mathematical errors.
The only mathematical error in the set passed verification because it lay in the translation of the
problem statement into the symbolic formalization, which confirms the phenomenon of silent errors and
delimits the reach of verification. Following this result the system now tells the teacher precisely
what was checked in each question. The educational product comprises the system, a curated bank of 85
questions and a teacher's manual, all freely available.

\medskip
\noindent\textbf{Keywords:} Mathematics Education. Artificial intelligence. Language models.
Multi-agent systems. Symbolic verification. Question authoring.
